\documentclass[11pt,a4paper]{article}

\usepackage[T1]{fontenc}
\usepackage[utf8]{inputenc}
\usepackage{lmodern}
\usepackage{amsmath,amssymb}
\usepackage{booktabs}
\usepackage{geometry}
\usepackage{xcolor}
\usepackage{listings}

\geometry{margin=2.5cm}

\lstset{
    basicstyle=\ttfamily\small,
    columns=fullflexible,
    frame=single,
    breaklines=true,
    showstringspaces=false
}

\title{The Two-Tape Matching Method}
\author{}
\date{}

\begin{document}

\maketitle

\section{The Two-Tape Matching Method}

Consider two one-dimensional tapes of identical length \(L\).

\subsection{Tape A: Buyers}

For each country \(c\), let
\[
A_c=\{b_1,b_2,\ldots,b_N\}
\]
be the ordered set of buyers belonging to that country.

The tape is divided into \(N\) equal segments, one for each buyer.
If the tape is normalized to unit length,
\[
L=1,
\]
the starting point of buyer \(i\) is
\[
x_i=\frac{i}{N},
\qquad i=0,\ldots,N-1.
\]

Each buyer therefore occupies an interval of length
\[
\Delta_A=\frac{1}{N}.
\]

The buyer stores only its own identifier, \texttt{buyer.uid}.

\subsection{Tape B: Suppliers}

For each sector \(s\), let
\[
B_s=\{s_1,s_2,\ldots,s_M\}
\]
be the suppliers producing intermediate goods in that sector.

Each supplier is assigned a positive weight
\[
v_j=\texttt{targetAddedValue}_j.
\]

To favour domestic suppliers, the weight is modified according to the
buyer's country:
\[
w_j=
\begin{cases}
k\,v_j, & \text{if } \operatorname{country}_j=c,\\[4pt]
v_j, & \text{otherwise},
\end{cases}
\]
where
\[
k>1
\]
is the domestic-preference factor.

The suppliers are randomly permuted once:
\[
\pi:\{1,\ldots,M\}\longrightarrow\{1,\ldots,M\},
\]
producing the shuffled tape
\[
B_s^\pi.
\]

The cumulative weights are then computed:
\[
C_j=\sum_{h=1}^{j} w_{\pi(h)},
\]
with
\[
W=C_M=\sum_{j=1}^{M}w_j.
\]

Supplier \(\pi(j)\) occupies the interval
\[
[C_{j-1},C_j),
\qquad C_0=0,
\]
whose length is
\[
C_j-C_{j-1}=w_{\pi(j)}.
\]

Thus, suppliers with larger weights occupy longer portions of the tape.

\section{Matching the Two Tapes}

The buyer tape is rescaled to the same total length \(W\) as the supplier
tape. The position of buyer \(i\) becomes
\[
p_i=x_iW.
\]

The assigned supplier is the unique supplier at position \(j\) satisfying
\[
C_{j-1}\leq p_i<C_j.
\]

Geometrically, the buyer is projected vertically onto the supplier tape.
The supplier interval containing that point determines the buyer's initial
supplier.

No independent random draw is performed for each buyer.

\section{Efficient Interval Search with \texttt{numpy.searchsorted}}

The cumulative vector
\[
(C_1,C_2,\ldots,C_M)
\]
is non-decreasing and, when all weights are positive, strictly increasing.

Given the buyer positions
\[
(p_1,p_2,\ldots,p_N),
\]
NumPy performs the interval searches with

\begin{lstlisting}[language=Python]
positions = np.searchsorted(
    cumulative_weights,
    buyer_points,
    side="right",
)
\end{lstlisting}

For each buyer position \(p_i\), \texttt{searchsorted} returns the smallest
index \(j\) such that
\[
C_j>p_i.
\]

This is precisely the index of the supplier interval containing the buyer.

The computational cost is
\[
O(M)
\]
for constructing the cumulative weights and approximately
\[
O(N\log M)
\]
for locating all buyer positions by binary search.

The searches are executed inside NumPy's compiled implementation, avoiding
a Python loop over all suppliers for each buyer.

\subsection{Why \texttt{side="right"}?}

Supplier intervals are defined as
\[
[C_{j-1},C_j).
\]

If a buyer lies exactly on a boundary,
\[
p_i=C_j,
\]
it must be assigned to the following interval. The option
\texttt{side="right"} implements this convention.

It also skips zero-length intervals generated by suppliers whose weight is
zero.

\section{Systematic Placement}

The recommended option is systematic placement.

Instead of using the exact beginning of every buyer segment,
\[
x_i=\frac{i}{N},
\]
all positions are translated by the same random offset
\[
u\sim U(0,1),
\]
giving
\[
x_i=\frac{i+u}{N}.
\]

After rescaling to tape length \(W\),
\[
p_i=\frac{i+u}{N}W.
\]

Geometrically, the entire buyer tape is shifted by a random distance smaller
than one buyer segment before being projected onto the supplier tape.

Only one random number is generated for the whole country-sector matching.

\section{Alternative Placement Schemes}

\subsection{Starts: deterministic placement}

The starting point of each buyer segment is used:
\[
x_i=\frac{i}{N}.
\]

Once the supplier tape has been shuffled, the matching is entirely
deterministic. No additional random numbers are generated.

\subsection{Systematic placement}

A single common random offset is used:
\[
x_i=\frac{i+u}{N},
\qquad u\sim U(0,1).
\]

Its main properties are:

\begin{itemize}
    \item only one random number is required for the entire tape;
    \item buyer positions remain perfectly equally spaced;
    \item supplier weights are represented very accurately;
    \item sampling variability is low;
    \item the geometric interpretation of the two aligned tapes is preserved.
\end{itemize}

\subsection{Stratified placement}

Each buyer receives an independent random position inside its own segment:
\[
x_i=\frac{i+u_i}{N},
\qquad u_i\sim U(0,1).
\]

Its main properties are:

\begin{itemize}
    \item one random number is generated for each buyer;
    \item every buyer remains inside its own segment;
    \item buyer positions are no longer equally spaced;
    \item stochastic variability is higher than under systematic placement.
\end{itemize}

\section{Comparison}

\begin{table}[htbp]
\centering
\begin{tabular}{lccc}
\toprule
Method & Random numbers & Buyer spacing & Variability \\
\midrule
Starts      & \(0\)   & perfectly regular & none after shuffling \\
Systematic  & \(1\)   & perfectly regular & low \\
Stratified  & \(N\)   & irregular          & higher \\
\bottomrule
\end{tabular}
\caption{Comparison of buyer-placement methods.}
\end{table}

The systematic method provides the best compromise between computational
efficiency, proportional representation of supplier weights, and the
geometric interpretation of the two aligned tapes.

\end{document}
